% ===================================================================
%  அட்டவணை எண் 14 — தெரு. வணிகர்கள்.
%
%  Traduction, faite en 2026, en tamoul d'aujourd'hui, sur l'ido de
%  texto/io. Regles de la colonne : voir 10-attavanai-01.tex. Recit a
%  la premiere personne, sans scenes. Ecrit avec outils/ordo.py sous
%  les yeux.
%
%  SEIZE PAIRES CONTRE TRENTE-TROIS AU TABLEAU 13, QUI COMPTE ONZE
%  RENVOIS DE MOINS. C'est la quatrieme colonne du releve a mesurer
%  cet ecart sur ces deux tableaux : le marathe avait releve neuf
%  inversions contre trois, le telougou dix-neuf contre neuf, le
%  coreen trente-trois paires contre seize, et le tamoul retrouve
%  exactement les memes trente-trois contre seize — puisque
%  parigi.py ne lit que l'ido et sort la meme liste pour toutes les
%  colonnes. Le 13 attache, le 14 enumere : le marchand de
%  parapluies, le ferblantier, l'armurier, les fiacres, les voitures
%  de maitre — et deux renvois joints par « e » ne sont pas en
%  rapport de dependance.
%
%  LES METIERS DE CASTE, TROISIEME FOIS DANS CETTE COLONNE. Le
%  cordonnier du tableau 7 avait deja pose la question ; ici
%  reviennent le coiffeur (23) et le tailleur (62). Le tamoul a des
%  noms de caste pour ces deux emplois, et ils servent d'injures. Le
%  coiffeur s'ecrit முடி திருத்துபவர், comme au tableau 7 ; le
%  tailleur s'ecrit தையல்காரர், qui est le mot ordinaire des
%  enseignes et ne designe personne. On ecrit la fonction, et on ne
%  la deguise pas non plus : தையல்காரர் est ce qui est peint sur les
%  boutiques de Chennai en 2026.
%
%  QUATRE-VINGT-DIX-NEUF RENVOIS, ET PRESQUE TOUS SONT DES METIERS.
%  C'est le tableau ou une langue montre le mieux ce qu'elle a
%  refait et ce qu'elle a garde. Le tamoul a ses mots pour tout ce
%  qui existait avant l'imprimerie : தையல்காரர், கொல்லர்,
%  தொப்பிக்காரர், நகைக்கடைக்காரர், பூக்காரி, துப்புரவுத் தொழிலாளர்.
%  Il emprunte pour ce qui est arrive avec la machine :
%  ஸ்டெனோகிராஃபர், தட்டச்சர், லித்தோகிராஃபர். La frontiere passe a
%  la meme place qu'au tableau 1 — le mobilier scolaire d'un cote,
%  les sciences de l'autre — et c'est la troisieme fois que ce
%  releve la voit tomber la.
%
%  DEUX PAIRES DE CE TABLEAU ONT ELARGI parigi.py, ET C'EST ECRIT
%  DANS LISEZ-MOI § 9 : « tamburestro (42) sequata da la tamburisti
%  (43) » et « porto-triciklo (80) duktata da grumo (79) » sont des
%  participes PASSIFS, que le motif ne prenait pas. Le livret est
%  passe de 386 a 408 paires a cause de ces deux-la. Elles sont donc
%  ecrites ici avec une attention particuliere : c'est le tableau qui
%  a paye pour l'outil.
% ===================================================================

%%K t14-apar-1 apar
\VUsaut{4.0mm}
\VUtitre{15.0pt}{60}{அட்டவணை எண் 14}
\VUsaut{3.0mm}

%%K t14-tit-1 sub
\VUpk{11.4pt}{40}{தெரு. --- வணிகர்கள்.}
\VUsaut{3.0mm}

%%K t14-01-1 p
1. --- நாட்டுப்புறத்தில் வாழும் என் நண்பன் இயோவான்னெஸ் என் குடும்பத்தில்
மூன்று நாள்களைக் கழிக்க வந்தான். \VUgras{அது வரை} அவனுக்குப் பெரிய
நகரத்தைப் பார்க்கும் வாய்ப்பு கிடைத்திருக்கவில்லை ; அதனால் நம் நாள்கள்
அனைத்தையும் உலாவிக் கழிக்கிறோம் ; அவனுக்குத் தெரியாதவற்றை நான்
விளக்குகிறேன்.

%%K t14-02-1 p
2. --- முதலில் \VUgras{வான் ஆய்வகத்தைப்} \textsuperscript{(1)} பார்க்கப்
போனோம் ; அது அவனுக்கு மிகவும் சுவாரசியமாக இருந்தது. திரும்பி வரும்
போது, \VUgras{தெரு} \textsuperscript{(3)} வழியாக வந்தோம் — அங்கே
\VUgras{துருக்கிய குளியல் இல்லம்} \textsuperscript{(2)} இருக்கிறது ; அங்கே ஒரு
\VUgras{இறுதிச் சடங்கைச்} \textsuperscript{(4)} சந்தித்தோம். \VUgras{இறுதி
ஊர்வலத்திற்கு} முன்னால் \VUgras{குருக்கள் குழு} நடந்தது ; அது
\VUgras{பிணவண்டிக்கு} முன்னால் சென்றது ; அந்த வண்டியில்
\VUgras{சவப்பெட்டி} வைக்கப்பட்டிருந்தது. பல நண்பர்கள் \VUgras{துக்கம்
அனுசரிக்கும்} குடும்பத்துடன் வந்தார்கள்.

%%K t14-02-2 p
ஊர்வலம் கடந்த பிறகு, பெரிய \VUgras{பீர் விடுதிக்கு}
\textsuperscript{(5)} முன்னால் தெருவைக் கடந்தோம் ; அங்கே பல
\VUgras{வாடிக்கையாளர்கள்} மொட்டை மாடத்தில் \VUgras{பீர்}
குடித்தார்கள் ; கண்ணாடி போட்ட \VUgras{மேற்கூரை} \textsuperscript{(6)}
அவர்களுக்கு நிழல் தந்தது.

%%K t14-03-1 p
3. --- \VUgras{மது வணிகரின்} \textsuperscript{(7)} கடைக்கும்
\VUgras{இசைக் கருவி வணிகரின்} \textsuperscript{(10)} கடைக்கும் இடையே
வைத்திருந்த \VUgras{சின்னம்} இயோவான்னெஸை மிகவும் வியப்பில்
ஆழ்த்தியது. அதன் பொருள் என்ன என்று என்னிடம் கேட்டான் ; அது ஆங்கிலேய
\VUgras{தூதரகத்தைக்} \textsuperscript{(8)} குறிக்கிறது என்று பதில்
சொல்லி, பால்கனியில் இருக்கும் \VUgras{தூதரையும்}
\textsuperscript{(9)} அவனுக்குக் காட்டினேன்.

%%K t14-tit-2 sub
\VUpk{10.2pt}{40}{கடைகளைப் பார்த்தல்.}
\VUsaut{3.0mm}

%%K t14-04-1 p
4. --- இயல்பாகவே \VUgras{இசைக் கருவிகளின்} காட்சிக்கு முன்னால்
நின்றோம் — \VUgras{ஹார்மோனியங்கள்}, \VUgras{ஆர்கன்கள்},
\VUgras{பியானோக்கள்} ; \VUgras{இசைப் பாடல்களின்} படங்கள் போட்ட
அட்டைகளையும் பியானோவுக்கான \VUgras{இசைத் தாள்களையும்} பார்த்தோம் ;
பெயர்ப் பலகையாகப் பயன்படும் பொன் முலாம் பூசிய மரத்தால் ஆன
\VUgras{யாழையும்} வியந்தோம்.

%%K t14-05-1 p
5. --- \VUgras{பெருக்குபவர்கள்} \textsuperscript{(11)} உண்டாக்கிய
தூசி மேகங்களும் \VUgras{குப்பை வண்டியும்} \textsuperscript{(12)}
\VUgras{தளவாட வணிகரின்} \textsuperscript{(13)} அழகான
\VUgras{தளவாடங்களுக்கு} முன்னாலும், \VUgras{தகரக் கடைக்காரரின்}
\textsuperscript{(14)} \VUgras{அடுப்புகள்} \VUgras{விளக்குகள்}
காட்சிக்கு முன்னாலும் விரைவாகக் கடந்து போகும்படி நம்மைக்
கட்டாயப்படுத்தின.

%%K t14-06-1 p
6. --- மேலும் இந்த இடத்தில் நடைபாதையை விட்டு இறங்க
வேண்டியிருந்தது ; ஏனெனில் அது மூட்டைகளால் அடைபட்டிருந்தது ; அவற்றை
ஒரு \VUgras{போக்குவரத்து வண்டியிலிருந்து} \textsuperscript{(16)}
இறக்கி, இப்போதுதான் வாடகைக்கு எடுத்த ஒரு \VUgras{கிடங்கில்}
\textsuperscript{(15)} வைக்கிறார்கள்.

%%K t14-06-2 p
அதனால் \VUgras{வண்டி செய்பவரின்} \textsuperscript{(17)} அழகான
வண்டிகளைப் பார்க்க முடியவில்லை ; ஆனால் \VUgras{மூட்டை கட்டுபவர்}
\textsuperscript{(19)} ஒரு பெட்டி செய்வதைப் பார்க்க நின்றோம் ; அதை
அவர் தன் அயலவருக்குச் செய்தார் — \VUgras{மிதிவண்டி, மகிழுந்து
வணிகருக்கு} \textsuperscript{(18)}. பிறகு நேரம் கொஞ்சம் ஆகி
விட்டதால் வீட்டுக்குத் திரும்பினோம்.

%%K t14-tit-3 sub
\VUpk{10.2pt}{40}{நகரத்தில் உலா.}
\VUsaut{3.0mm}

%%K t14-07-1 p
7. --- மறு நாள் என் தாய் இயோவான்னெஸின் படம் எடுக்கும்படி
முன்மொழிந்தாள் ; பிறகு \VUgras{ஒளிப்படக் கலைஞரிடம்}
\textsuperscript{(20)} அவனை அழைத்துப் போகும் பணியை எனக்குத் தந்தாள்.
மதிய உணவுக்குப் பிறகு அந்தக் கலைஞரின் \VUgras{பணிமனைக்குப்}
போனோம் ; அங்கே நெடு நேரம் காத்திருக்க வேண்டியிருந்தது. நேரத்தைக்
கழிக்கச் சுவரில் தொங்கும் \VUgras{உருவப்படங்களைப்}
\textsuperscript{(22)} பார்த்தோம்.

%%K t14-07-2 p
நம் முறை வந்த போது வேலை நெடு நேரம் நீடிக்கவில்லை ; என் நண்பன்
\VUgras{ஒளிப்படக் கருவிக்கு} \textsuperscript{(21)} முன்னால் நிற்பதை
முடித்த உடனேயே கீழே இறங்கினோம்.

%%K t14-08-1 p
8. --- முதல் தளத்தில் \VUgras{முடி திருத்துபவரின்}
\textsuperscript{(23)} கதவு திறந்திருந்ததால், அதன் வழியாக அவருடைய
உதவியாள் ஒரு வாடிக்கையாளரின் தலை முடியை வெட்டுவதைப் பார்த்தோம்.

%%K t14-08-2 p
தெருவுக்கு வந்த பிறகு \VUgras{புகையிலைக் கடைக்கு}
\textsuperscript{(24)} முன்னால் கடந்தோம். சிவப்பு \VUgras{லாந்தரும்}
\textsuperscript{(25)}, \VUgras{பெயர்ப் பலகையாகப்} \textsuperscript{(26)}
பயன்படும் \VUgras{பருமனான புகைக் குழலும்} இயோவான்னெஸை அவ்வளவு
மகிழ்வித்ததால், \VUgras{நசியப் பொடி முகர்ந்து} \textsuperscript{(27)}
பேசிக் கொண்டிருந்த இரண்டு முதியவர்களுடன் அவன் மோதி விட்டான்.

%%K t14-tit-4 sub
\VUpk{10.2pt}{40}{இனிப்புக் கடை.}
\VUsaut{3.0mm}

%%K t14-09-1 p
9. --- \VUgras{நாணய மாற்றுபவரின்} \textsuperscript{(29)}
\VUgras{கண்ணாடிப் பேழையில்} காட்சிக்கு வைத்திருந்த எல்லா நாட்டு
\VUgras{வங்கித் தாள்களும்} \VUgras{நாணயங்களும்} சில கணங்கள் நம்
கவனத்தை ஈர்த்தன ; ஆனால் சிற்றுண்டி நேரம் ஆனதால், மேலும் நிற்காமல்
\VUgras{இனிப்புக் கடைக்காரரிடம்} \textsuperscript{(31)} நுழைந்தோம்.

%%K t14-10-1 p
10. --- கடையில் இருந்த எல்லா \VUgras{சுவையான இனிப்புகளையும்}
பார்த்த போது — \VUgras{கேக்குகள்}, \VUgras{தேன் கேக்குகள்},
\VUgras{பிஸ்கட்டுகள்}, எல்லா வகை \VUgras{மிட்டாய்கள்} — எளிதாகத்
தேர்ந்தெடுக்க முடியவில்லை. கடைசியாக இயோவான்னெஸ் \VUgras{ஏடு
கேக்கைத்} தேர்ந்தெடுத்தான் ; நான் சிறிய \VUgras{செர்ரிப் பணியாரம்}
மட்டுமே எடுத்தேன் ; ஏனெனில் இனிப்புகள் \VUgras{என் பற்களை வலிக்க
வைக்கின்றன} ; \VUgras{பல் மருத்துவரிடம்} \textsuperscript{(30)}
அடிக்கடி போக நான் விரும்பவே இல்லை.

%%K t14-tit-5 sub
\VUpk{10.2pt}{40}{தட்டச்சு.}
\VUsaut{3.0mm}

%%K t14-11-1 p
11. --- என் நண்பன் கேள்விப்பட்டிருந்த ஆனால் பார்த்திராத
\VUgras{தட்டச்சு இயந்திரத்தை} அவனுக்குக் காட்ட, ஓர் \VUgras{எழுது அறைக்குள்} \textsuperscript{(32)} நுழைந்தோம் ;
அது என் \VUgras{ஒன்றுவிட்ட சகோதரரினுடையது} \textsuperscript{(34)}.
அவர் தன் கடிதங்களைத் தன்
\VUgras{செயலாளருக்குச்} \textsuperscript{(35)} சொல்லிக்
கொண்டிருந்தார் ; அந்தச் செயலாளர் \VUgras{ஸ்டெனோகிராஃபர்}. கடிதம்
முடிந்த உடனேயே \VUgras{தட்டச்சர்} \textsuperscript{(33)} அதைத் தன்
இயந்திரத்தில் நகல் எடுத்தார். என் உறவினரின் வேலைகளுக்கு இடையூறு
செய்ய விரும்பாமல், சில கணங்கள் மட்டுமே அவரிடம் இருந்தோம்.

%%K t14-12-1 p
12. --- என் ஒன்றுவிட்ட சகோதரரின் அலுவலகத்திற்குக் கீழே பெரிய
\VUgras{கடிகார நகைக் கடைக்காரர்} \textsuperscript{(37)}
குடியிருக்கிறார் ; அவர் தங்க வேலைக்காரரும் கூட. அவருடைய சிறந்த
கடையில் பல \VUgras{கடிகாரங்கள்}, \VUgras{ஊசல் கடிகாரங்கள்},
\VUgras{சட்டைப் பை கடிகாரங்கள்}, \VUgras{சிறு சங்கிலிகள்},
விலையுயர்ந்த \VUgras{நகைகள்} இருக்கின்றன — எடுத்துக்காட்டாக
\VUgras{முத்து மாலைகள்}, தங்க \VUgras{கை வளையல்கள்},
\VUgras{காதணிகள்}, \VUgras{இரத்தினக் கற்களாலும்}
\VUgras{வைரங்களாலும்} அலங்கரித்த \VUgras{மோதிரங்கள்}. கண்ணாடிப்
பேழைகளுள் ஒன்று \VUgras{தங்கப் பொருள்களுக்கே} ஒதுக்கப்
பட்டிருக்கிறது ; அதில் வெள்ளி \VUgras{உணவுக் கரண்டிகளின்}
முக்கியமான தொகுப்பு இருக்கிறது.

%%K t14-13-1 p
13. --- தெருவின் மூலையில் திரும்பும் போது, வீட்டின்
\VUgras{எண்ணுக்கு} அருகில் வைத்த \VUgras{பலகையை}
\textsuperscript{(36)} மாற்றியிருப்பதைக் கவனித்தேன். என் கவனிப்பை
உரக்கச் சொல்லப் போனேன் ; அப்போது \VUgras{படை இசையின்}
மகிழ்ச்சியான ஒலிகள் கேட்டன. படைப் பிரிவை நோக்கி ஓடினோம் ;
\VUgras{அது} \VUgras{தொப்பிக்காரரின்} \textsuperscript{(39)},
\VUgras{கையுறை வணிகரின்} \textsuperscript{(40)} கடைகளுக்கு முன்னால்
கடக்கப் போகிறது என்பதையும், \VUgras{மூக்குக் கண்ணாடிகள்},
\VUgras{காதுக் கண்ணாடிகள்}, \VUgras{சிறு தொலைநோக்கிகள்} நிறைந்த
\VUgras{கண்ணாடி வணிகரின்} \textsuperscript{(38)} கண்ணாடிப் பேழைக்கு
முன்னால் நின்றாலும் அதன் அணிவகுப்பை அதே அளவு நன்றாகப் பார்க்க
முடியும் என்பதையும் எண்ணிப் பார்க்காமல்.

%%K t14-tit-6 sub
\VUpk{10.2pt}{40}{அணிவகுப்பு.}
\VUsaut{3.0mm}

%%K t14-14-1 p
14. --- \VUgras{படைப் பிரிவுக்கு} \textsuperscript{(41)} முன்னால்
\VUgras{முரசுத் தலைவர்} \textsuperscript{(42)} நடந்தார் ; அவரைப்
பின்தொடர்ந்து \VUgras{முரசு அடிப்பவர்களும்} \textsuperscript{(43)},
\VUgras{எக்காளக் காரர்களும்} \textsuperscript{(44)}, \VUgras{இசைக்
குழுவினர்} முழுவதும் வந்தார்கள். \VUgras{தளபதி}
\textsuperscript{(45)} தன் உதவி அதிகாரியுடன் சதுக்கத்தில் இருந்தார் ;
\VUgras{அணிவகுப்பைப்} பார்ப்பதற்காக \VUgras{நீர்ப் பீச்சிக்கு}
\textsuperscript{(49)} அருகில் நின்றிருந்தார் — \VUgras{நீரூற்றின்}
\textsuperscript{(48)} பீச்சிக்கு அருகில்.

%%K t14-14-2 p
\VUgras{தலைமைப் பணி அதிகாரிகள்} \textsuperscript{(46)},
\VUgras{படைத் தலைவர்}, \VUgras{துணைப் படைத் தலைவர்} ஆகியோர்
வாளால் அவரை வணங்கினார்கள். \VUgras{பிரிவுத் தலைவர்கள்} தங்கள்
\VUgras{பிரிவுகளுக்கு} \textsuperscript{(47)} முன்னால் இருந்தார்கள் ;
ஒவ்வொரு \VUgras{தலைவரும்} தன் \VUgras{படைக் குழுவுக்குக்} கட்டளை
இட்டார் ; அதே வேளையில் \VUgras{துணை அதிகாரிகளும்},
\VUgras{இளநிலை அதிகாரிகளும்}, \VUgras{கீழ் அதிகாரிகளும்} தங்கள்
வீரர்களுக்கு அருகில் வழக்கமான இடத்தைப் பிடித்தார்கள்.

%%K t14-15-1 p
15. --- \VUgras{சந்திப்பில்} \textsuperscript{(50)} பல வழிப்போக்கர்கள்
கூட்டமாகக் கூடினார்கள் ; வணிகர்கள், \VUgras{குடை செய்பவர்}
\textsuperscript{(51)}, \VUgras{தகரக் கம்பி வேலைக்காரர்}
\textsuperscript{(53)}, \VUgras{ஆயுத வணிகர்} \textsuperscript{(54)}
\VUgras{வீரர்களைப்} பார்க்க வெளியே வந்தார்கள். எல்லா வண்டிகளும் —
\VUgras{வாடகை வண்டிகள்} \textsuperscript{(52)}, \VUgras{தனி வண்டிகள்}
\textsuperscript{(57)}, \VUgras{நீர்த் தெளிக்கும் பீப்பாய் வண்டியும்}
\textsuperscript{(56)} — நடைபாதை ஓரமாக நின்றன.

%%K t14-tit-7 sub
\VUpk{10.2pt}{40}{தெருவில் காட்சிகள்.}
\VUsaut{3.0mm}

%%K t14-15-2 p
ஒரு \VUgras{வணிகப் பயணி} \textsuperscript{(55)} கையில் தன்
\VUgras{மாதிரிப் பேழைகளைச்} சுமந்து விரைவாகத் தெருவைக் கடந்தார் ;
\VUgras{மேல் தளத்தில்} \textsuperscript{(59)} —
\VUgras{பொது வண்டியின்} \textsuperscript{(58)} மேல் தளத்தில் —
அமர்ந்திருந்தவர்கள் அந்தக் காட்சியை அனுபவிக்கத்
திரும்பினார்கள். பரிதாபமான \VUgras{ஊர் சுற்றும் பாடகர்} \textsuperscript{(60)} —
அவர் தன் \VUgras{சுழல் ஆர்கனுடன்} \textsuperscript{(61)} வந்தார் —
தன் \VUgras{பாடலை} நிறுத்த வேண்டியிருந்தது ; ஏனெனில் முரசுகளின்
ஆரவாரம் அவருடைய குரலை மூடி விட்டது.

%%K t14-16-1 p
16. --- படைப் பிரிவு \VUgras{துணிக் கடைக்கு} \textsuperscript{(64)}
முன்னால் வந்த போது, \VUgras{தையல் பெண்களும்} \textsuperscript{(63)}
\VUgras{தையல்காரர்களும்} \textsuperscript{(62)} சாளரத்தின் வழியாகப்
பார்க்கத் தங்கள் \VUgras{தையல் இயந்திரங்களையும்} வேலையையும்
விட்டு விட்டார்கள். \VUgras{வணிகர்} \textsuperscript{(65)} —
அவர் இப்போதுதான் புதிய \VUgras{உடைகளை} \textsuperscript{(66)} தன்
\VUgras{காட்சி இடத்தில்} வைத்திருந்தார் — \VUgras{கம்பளித் துணிகளையும்} \textsuperscript{(67)} \VUgras{நூல்
துணிகளையும்} உள்ளே வைக்க — அவை நடைபாதையில் \VUgras{கழிவு நீர்
வாய்க்கு} \textsuperscript{(68)} அருகில் வைக்கப்பட்டிருந்தன —
வேண்டியிருந்தது ; கூட்டம் அவற்றைக் கவிழ்த்து விடும் என்று அஞ்சி.
பழைய \VUgras{சுமந்து விற்பவர்}
\textsuperscript{(69)}கூட — அவரிடம் ஒரு காப்பாளர் காலுறைகளுக்கு
விலை பேசினாள் — தன் விற்பனையை முடிக்க ஒரு நடைபாதைக்குள் நுழைய
வேண்டியிருந்தது.

%%K t14-16-2 p
படைப் பிரிவைப் படை முகாம் வரை தொடர்ந்தோம் ; \VUgras{மேலும்} நம்
நாளைப் பற்றி மிகவும் மனநிறைவு அடைந்து இரவு உணவின் நேரத்தில்
திரும்பினோம்.

%%K t14-tit-8 sub
\VUpk{10.2pt}{40}{பல வணிகங்களும் தொழில்களும்.}
\VUsaut{3.0mm}

%%K t14-17-1 p
17. --- மறு நாள் என் தந்தை அங்குள்ள செய்தித்தாளின் அச்சகத்தைப்
பார்க்க முன்மொழிந்தார். ஆர்வத்துடன் ஏற்றுக் கொண்டோம்.

%%K t14-17-2 p
\VUgras{அச்சிடுபவர்} \VUgras{புத்தக வணிகரும்} \textsuperscript{(70)}
\VUgras{(= புத்தகம் விற்பவர்)}, \VUgras{பதிப்பாளரும்},
\VUgras{காகித வணிகரும்}, \VUgras{புத்தகம் கட்டுபவரும்} கூட. முதல்
தளத்தில் செய்தித்தாளின் \VUgras{ஆசிரியர் அறையையும்}
\textsuperscript{(71)} அமைத்திருக்கிறார் ; \VUgras{பொறிக்கும்
அறையையும்} — அங்கே \VUgras{கல்லச்சுக் கலைஞர்களும்}
\textsuperscript{(72)} \VUgras{செம்புப் பொறிப்பாளர்களும்}
\textsuperscript{(73)} வேலை செய்கிறார்கள் — கடைசியாக
\VUgras{அச்சு அடுக்கும் அறையையும்} ; அங்கே \VUgras{அச்சுத்
தொழிலாளர்கள்} \textsuperscript{(74)} இருக்கிறார்கள். உண்மையான
\VUgras{அச்சகமும்} \textsuperscript{(75)}, \VUgras{அச்சு
இயந்திரங்களும்}, பல இயந்திரங்களும் தரைத் தளத்தில் இருக்கின்றன.

%%K t14-tit-9 sub
\VUpk{10.2pt}{40}{இயோவான்னெஸ் நிறையக் கேள்விகள் கேட்கிறான்.}
\VUsaut{3.0mm}

%%K t14-18-1 p
18. --- அச்சகத்தை விட்டு வெளியே வரும் போது இயோவான்னெஸ் மிகவும்
வியந்து நின்றான் ; \VUgras{சிறு சாமான் கடைக்கு} \textsuperscript{(77)}
எதிரே ஒரு தூணின் மேல் வைத்திருந்த சிறிய கண்ணாடிப் பேழைக்கு
முன்னால் ; அது எதற்குப் பயன்படுகிறது என்று கேட்டான். அது \VUgras{தீ
எச்சரிக்கைப் பெட்டி} \textsuperscript{(76)} என்று என் தந்தை அவனுக்கு
விளக்கினார். மேலும் நகரத்தில் எல்லாமே என் நண்பனுக்கு வியப்பாக
இருந்தது — எளிய \VUgras{கல் நீரூற்று} \textsuperscript{(78)}
முதல், இலேசான \VUgras{சுமை முச்சக்கர வண்டி} \textsuperscript{(80)} வரை —
அதைச் சீருடை அணிந்த \VUgras{சிறு பணியாள்} \textsuperscript{(79)}
ஓட்டுகிறான் —, \VUgras{அஞ்சல் வண்டி} \textsuperscript{(81)} வரை.

%%K t14-19-1 p
19. --- என் தந்தை சில கணங்கள் நம்மை விட்டு \VUgras{வரி
வசூலிப்பவருடன்} \textsuperscript{(83)} பேசப் போன போது — அவர் ஒரு
\VUgras{வழக்கறிஞருடனும்} \textsuperscript{(82)} ஒரு
\VUgras{நோட்டரியுடனும்} \textsuperscript{(84)} பேச நின்றிருந்தார் —
தெருவின் நடமாட்டத்தைக் கண்களால் பின்தொடர்ந்து மகிழ்ந்தோம். மிகப்
பலவகை மனிதர்கள் நம் முன்னால் கடந்தார்கள் : ஒரு \VUgras{இனிப்புக்
கடைப் பணியாளன்} \textsuperscript{(85)} கூடையில் சிறந்த
\VUgras{பணியாரம்} சுமந்து ஒரு \VUgras{ஆடை வணிகருக்குப்}
\textsuperscript{(86)} பின்னால் வந்தான் ; ஒரு \VUgras{லாந்தர்
ஏற்றுபவர்} \textsuperscript{(87)} ஒரு \VUgras{வாயு ஊழியருடன்}
\textsuperscript{(88)} பேசிக் கொண்டிருந்தார் ; ஒரு \VUgras{கன்னியாஸ்திரி}
\textsuperscript{(89)} கையில் ஒரு அனாதைச் சிறுமியைப் பிடித்துக்
கொண்டு தன் மடத்திற்குத் திரும்பிக் கொண்டிருந்தாள்.

%%K t14-tit-10 sub
\VUpk{10.2pt}{40}{வீடு திரும்பும் போது.}
\VUsaut{3.0mm}

%%K t14-20-1 p
20. --- என் தந்தை நம்மைச் சேர்ந்த கணத்தில், இயோவான்னெஸ் வெடித்துச்
சிரித்தான் ; முற்றிலும் கரி படிந்த இரண்டு \VUgras{புகைபோக்கி
சுரண்டுபவர்களின்} \textsuperscript{(91)} அழுக்கான முகத்தையும்
வித்தியாசமான உடையையும் பார்த்து.

%%K t14-21-1 p
21. --- \VUgras{பூக்காரியிடம்} \textsuperscript{(92)} என் தாய்க்கு ஒரு
செடி வாங்க விரும்பினேன் ; ஆனால் அதைச் சுமந்து போவது மிகவும்
பாரமாக இருக்கும் என்றும், \VUgras{ஆரஞ்சுப் பழங்கள்} வாங்குவது
மேலானது என்றும் என் தந்தை எனக்குக் காட்டினார்.
\VUgras{பழக்காரியை} \textsuperscript{(94)} அணுகினோம் ;
\VUgras{ஆசிரியை} \textsuperscript{(93)} — அவள் தன் மாணவிக்காகத்
தேர்ந்தெடுத்துக் கொண்டிருந்தாள் — தன் கொள்முதலை முடித்த பிறகு
நமக்குப் பரிமாறச் செய்தோம்.

%%K t14-22-1 p
22. --- வீட்டுக்குத் திரும்பும் போது பல அறிமுகமானவர்களைச்
சந்தித்தோம் : என் குடும்பத்தின் பழைய தோழி — அவள் தன்
\VUgras{துணைப் பெண்மணியின்} \textsuperscript{(95)} கையைப்
பிடித்திருந்தாள் —, பாலங்களுக்கும் சாலைகளுக்கும் உரிய
\VUgras{பொறியாளர்} \textsuperscript{(96)}, என் ஒன்றுவிட்ட
சகோதரிகளுள் ஒருத்தி தன் \VUgras{குழந்தை காப்பாளருடன்}
\textsuperscript{(98)}.

%%K t14-22-2 p
பிறகு என் தாயின் \VUgras{தொப்பி செய்பவளையும்} \textsuperscript{(97)},
நகரத்திற்குப் புதிய இரண்டு \VUgras{துறவிகளையும்}
\textsuperscript{(99)} சந்தித்தோம் ; அவர்கள் நிலையத்திற்கான வழியை
என் தந்தையிடம் கேட்க நம்மிடம் வந்தார்கள்.

\VUsaut{6.0mm}
\VUfilet{22mm}
